%------------------------------------------------- Chapter 4: implementation
\chapterseparator{4}{Implementation, Deployment, and Agent Evaluation}

\chapter{Implementation, Deployment, and Agent Evaluation}

\section{Introduction}

This chapter details the technical implementation, containerised deployment, and comprehensive empirical evaluation of the AI-powered recommendation system. Following the CRISP-DM engineering lifecycle, the chapter is structured into two main parts. The first part describes the implementation of the core components: the layered system architecture, the automated extraction pipeline, the deterministic scoring engine, the asynchronous FastAPI microservice, the WordPress shortcode integration, the interactive web interfaces, and the multi-container Docker Compose deployment. The second part presents a formal evaluation protocol based on 22 realistic customer scenarios, analysing prompt sensitivity, key performance metrics, failure modes, and positioning relative to industrial commercial agents.

\section{System architecture and layered design}

The system is constructed as a decoupled, multi-tiered architecture where each layer communicates with the layer immediately below it via strictly typed interfaces. Figure~\ref{fig:sys_arch} illustrates the complete architectural blueprint and data flow.

\placeholderscreenshot{system_architecture.png}{Detailed multi-layer architecture and data flow of the recommendation system}{fig:sys_arch}

A core design virtue of this architecture is graceful degradation:
\begin{itemize}
  \item The deterministic scoring engine operates independently as an offline library;
  \item The FastAPI microservice serves questionnaire choices and computed routines without invoking any external language model;
  \item The generative conversational agent acts as an autonomous top-level controller, utilizing the lower deterministic tiers as specialized tools. If the third-party LLM provider experiences an outage, the core e-commerce recommendation service continues to operate uninterrupted.
\end{itemize}

\section{The automated extraction pipeline}

\subsection{Database ingestion and retention criteria}

Products are ingested from the MariaDB database via an optimized SQL extraction script (\texttt{pipeline/1\_export.sql}) that joins \texttt{posts}, \texttt{postmeta}, and taxonomy tables while filtering out merchandising buckets (such as \emph{Promotions} and \emph{Accueil}).

The ingestion retention rule proved essential: an initial filter required products to have a non-empty \texttt{post\_content} long description, which silently discarded 37 of the store's 52 commercial packs (which rely on image banners rather than body text). The retention condition was updated to accept items possessing either a title, a short excerpt, or a pack identifier.

\subsection{Multi-level pattern extraction}

Figure~\ref{fig:pipeline_flow} presents the execution flowchart of the extraction pipeline. Texts are stripped of HTML formatting, diacritics are normalized, and tokens are matched against prioritized regex patterns.

\placeholderscreenshot{extraction_pipeline_flow.png}{Flowchart of the multi-stage extraction, normalisation and pack resolution pipeline}{fig:pipeline_flow}

Table~\ref{tab:extraction} details the extraction performance and progressive improvements achieved between initial development (Version 1) and production delivery (Version 2).

\begin{table}[H]
\caption{Extraction pipeline performance and quality improvements across versions}
\label{tab:extraction}
\centering
\begin{tabularx}{\textwidth}{X r r}
\toprule
\textbf{Quality indicator} & \textbf{Initial version (v1)} & \textbf{Final version (v2)} \\
\midrule
Total catalogue products processed & 550 & 609 \\
Unclassified product format & 110 & \textbf{0} \\
Commercial packs successfully resolved & 12 & \textbf{52} \\
Topical cosmetic products with skin-type coverage & 70 & \textbf{110} \\
Products flagged for manual QA review & 49 & \textbf{5} \\
Full catalogue execution duration & $<$2\,s & $<$2\,s \\
\bottomrule
\end{tabularx}
\end{table}

\subsection{Pack resolution via TF-IDF rarity token matching}

To recover unlinked bundle compositions, the pipeline compares pack descriptions against non-pack titles. Table~\ref{tab:packs} shows that applying TF-IDF rarity weighting and threshold gates eliminated false-positive component associations entirely, securing 20 high-fidelity routine templates.

\begin{table}[H]
\caption{Pack composition recovery: comparison between unweighted and TF-IDF weighted matching}
\label{tab:packs}
\centering
\begin{tabularx}{\textwidth}{X r r}
\toprule
\textbf{Resolution metric} & \textbf{Unweighted matching} & \textbf{TF-IDF weighted matching} \\
\midrule
Packs with resolved contents & 39 / 52 & 20 / 52 \\
Constituent products accurately linked & 73 & 32 \\
Observed false-positive component linkages & Frequent & \textbf{Zero observed} \\
\bottomrule
\end{tabularx}
\end{table}

\section{The deterministic scoring engine}

\subsection{Multi-attribute scoring formulation}

Candidate items are scored by combining dermatological agreement, product popularity, and extraction confidence. Table~\ref{tab:scoring} outlines the explicit scoring terms.

\begin{table}[H]
\caption{Deterministic scoring weights and dermatological constraints}
\label{tab:scoring}
\centering
\begin{tabularx}{\textwidth}{L{5.4cm} X}
\toprule
\textbf{Evaluation criterion} & \textbf{Score adjustment or algorithmic action} \\
\midrule
Primary customer concern matched & $+15$ points \\
Secondary concern matched & $+15 - 2.5i$ points (floor $+4$) \\
No concern matched & \textbf{Strict exclusion from results} \\
Declared skin type matches & $+9$ points \\
Universal suitability (``all skin types'') & $+3$ points \\
Contradictory skin type declared & \textbf{Strict disqualification ($-\infty$)} \\
Undeclared skin type (unknown) & Eligible, 0 bonus points \\
Anatomical zone match (e.g. face) & $+5$ points \\
Adjacent zone off-target & $-4$ penalty points \\
Application timing match (morning/night) & $+4$ points \\
Verified customer rating & $+1.1 \times \text{rating}$ \\
Sales volume logarithm & $+0.7 \times \ln(1 + \text{sales})$ \\
Extraction confidence coefficient & Multiplied by $(0.6 + 0.4 \times \text{confidence})$ \\
Out-of-stock inventory status & $-6$ penalty points \\
\bottomrule
\end{tabularx}
\end{table}

\subsection{Routine assembly and conflict handling}

Topical skincare routines are constructed along five canonical steps: Cleanse, Tone, Treat (Serum), Moisturize, and Protect. A dedicated assembly routine handles ingestible supplements to prevent mixing topical creams with oral capsules. A dermatological conflict matrix prevents cross-recommending weight-loss teas for digestive discomfort.

\section{The API service and microservice engineering}

\subsection{RESTful endpoint specifications}

The backend engine is exposed via an asynchronous FastAPI service (\texttt{service/}). Table~\ref{tab:endpoints} summarizes the core endpoints.

\begin{table}[H]
\caption{FastAPI microservice endpoints and operations}
\label{tab:endpoints}
\centering
\begin{tabularx}{\textwidth}{L{2.0cm} L{3.0cm} X}
\toprule
\textbf{HTTP method} & \textbf{Endpoint route} & \textbf{Functional role} \\
\midrule
\texttt{GET}  & \texttt{/health} & Service health, catalogue size, snapshot freshness \\
\texttt{GET}  & \texttt{/options} & Available concerns and skin types with live product counts \\
\texttt{POST} & \texttt{/recommend} & Deterministic routine calculation \\
\texttt{POST} & \texttt{/interpret} & Natural-language query parsing into JSON profile \\
\texttt{POST} & \texttt{/chat} & ReAct conversational agent dialogue loop \\
\texttt{POST} & \texttt{/admin/refresh} & Token-secured hot catalogue re-extraction \\
\bottomrule
\end{tabularx}
\end{table}

Figure~\ref{fig:screenshot_fastapi_swagger} illustrates the interactive OpenAPI Swagger documentation interface provided automatically by FastAPI at \texttt{/docs}.

\placeholderscreenshot{screenshot_fastapi_swagger.png}{Interactive API documentation and endpoint testing via FastAPI Swagger UI}{fig:screenshot_fastapi_swagger}

\subsection{Two-speed catalogue refresh and security hardening}

Catalogue metadata is cached in-memory and periodically re-extracted in the background, whereas unit pricing and inventory stock counts are verified in real time from MariaDB for the returned items. The container executes under an unprivileged user, and database access is locked to a read-only \texttt{SELECT} role.

\section{Storefront integration and user interfaces}

\subsection{WordPress custom plugin}

A native WordPress plugin was developed (\texttt{scripts/make\_plugin.py}) providing two shortcodes: \texttt{[moteur\_recommandation]} for the interactive guided questionnaire, and \texttt{[conseiller\_beaute]} for the conversational agent dialogue. Requests are proxied server-side via PHP, hiding backend URLs and avoiding CORS issues.

Figure~\ref{fig:screenshot_wordpress_plugin} depicts the WordPress integration and shortcode deployment in the store's administration panel and front-end.

\placeholderscreenshot{screenshot_wordpress_plugin.png}{WordPress administration and WooCommerce store front-end displaying embedded recommendation shortcode}{fig:screenshot_wordpress_plugin}

\subsection{Interactive web demo interfaces}

A standalone web interface was built using responsive HTML5 and Vanilla JavaScript (\texttt{demo/index.html}). Figure~\ref{fig:screenshot_demo_questionnaire} shows the guided questionnaire interface with interactive concern chips and skin type filters. Figure~\ref{fig:screenshot_demo_routine} displays the generated 4-step routine with budget tracking and direct add-to-cart buttons.

\placeholderscreenshot{screenshot_demo_questionnaire.png}{Interactive recommendation web demo interface --- Guided questionnaire with concern chips and skin type selectors}{fig:screenshot_demo_questionnaire}

\placeholderscreenshot{screenshot_demo_routine.png}{Interactive recommendation web demo interface --- Generated 4-step personalized skincare routine with budget summary}{fig:screenshot_demo_routine}

\subsection{Conversational shopping assistant chat}

Figure~\ref{fig:screenshot_agent_chat} illustrates the conversational chat interface where users can converse naturally with the beauty assistant.

\placeholderscreenshot{screenshot_agent_chat.png}{Conversational shopping adviser chat interface with natural language inquiry and tool-grounded response cards}{fig:screenshot_agent_chat}

\section{Conversational agent and ReAct tool use}

\subsection{Agentic loop architecture}

Table~\ref{tab:agentvspipeline} contrasts a single-pass pipeline with the multi-turn ReAct agent. The model evaluates user utterances, queries tools, and iterates until sufficient information is gathered. Figure~\ref{fig:agent_cycle} diagrams the decision and execution cycle of the agent.

\begin{table}[H]
\caption{Architectural comparison between linear pipeline and ReAct agent}
\label{tab:agentvspipeline}
\centering
\begin{tabularx}{\textwidth}{L{4.6cm} X X}
\toprule
\textbf{System property} & \textbf{Linear pipeline} & \textbf{ReAct conversational agent} \\
\midrule
LLM invocations per request & Fixed (1) & Dynamic (1 to 8 iterations) \\
Next action decision & Hardcoded in logic & Decided dynamically by LLM \\
Conversational memory & None (stateless) & Full multi-turn dialog history \\
Available tool instrumentation & None & 3 dedicated deterministic tools \\
\bottomrule
\end{tabularx}
\end{table}

\placeholderscreenshot{agent_react_cycle.png}{Conversational agent decision cycle: user intent parsing, tool dispatching and response generation}{fig:agent_cycle}

The agent is instrumented with three tools:
\begin{itemize}
  \item \texttt{lister\_options}: Returns valid concerns and skin types present in the catalogue;
  \item \texttt{chercher\_produits}: Executes ranked search for targeted requests;
  \item \texttt{construire\_routine}: Assembles a multi-step routine respecting budget limits.
\end{itemize}

\subsection{Strict structural anti-hallucination guardrail}

To prevent hallucinations, product cards are rendered solely from verified database IDs returned by tool calls. The model cannot inject unverified items or falsified prices into the UI.

\section{Docker containerisation and deployment}

The complete microservice stack is containerised using Docker and Docker Compose. Figure~\ref{fig:screenshot_docker_deployment} shows the running container status and service logs.

\placeholderscreenshot{screenshot_docker_deployment.png}{Docker Compose multi-container deployment environment and running service status}{fig:screenshot_docker_deployment}

\section{Agent evaluation and analysis of results}

\subsection{Evaluation objectives and experimental protocol}

The evaluation evaluates whether the system delivers dermatologically safe, commercially viable, and conversationally coherent recommendations across diverse user profiles. A test suite of 22 realistic customer personas was formulated, covering common and rare concerns, conflicting requests, and edge cases.

\subsection{Typology of evaluation scenarios}

The 22 evaluation scenarios are categorised into three distinct typologies:
\begin{enumerate}
  \item \textbf{Standard single-concern queries}: Clear customer needs (e.g. teenage acne on oily skin, or anti-wrinkle care for mature skin);
  \item \textbf{Multi-concern and cross-domain requests}: Complex requests mixing dermatological issues with lifestyle constraints (e.g. pregnancy-induced hyperpigmentation combined with acute dehydration, or sleep disturbances paired with dry skin);
  \item \textbf{Adversarial and ambiguous edge cases}: Queries containing vague colloquialisms, contradictory skin types (e.g. claiming both severely dry and extremely oily skin), or requests for out-of-stock items.
\end{enumerate}

\subsection{Impact of prompt engineering on agent response quality}

Early testing revealed that minor phrasing variations in the system prompt dramatically influenced dialogue stability. Table~\ref{tab:prompt_impact} documents the observed failure modes and the prompt refinements implemented to resolve them.

\begin{table}[H]
\caption{Impact of prompt engineering refinements on agent conversational reliability}
\label{tab:prompt_impact}
\centering
\begin{tabularx}{\textwidth}{L{4.2cm} L{4.6cm} X}
\toprule
\textbf{Initial prompt issue} & \textbf{Observed failure behaviour} & \textbf{Engineered prompt correction} \\
\midrule
Unconstrained arithmetic & Model hallucinated inaccurate total routine prices & Forced total calculation to originate from tool return values only \\
Currency drift & Quoting prices in Euros rather than Tunisian Dinars & Injected explicit localized currency constraint (\texttt{DT}) into system header \\
Tool looping & Repeatedly calling \texttt{lister\_options} in an infinite loop & Implemented duplicate call detection and forced textual closure rule \\
Unsolicited cross-selling & Suggesting weight-loss teas for simple digestion queries & Introduced dermatological conflict boundary instructions \\
Card accumulation & Older recommendation cards persisted across dialogue turns & Frontend updated to replace stale cards with latest tool result set \\
\bottomrule
\end{tabularx}
\end{table}

\subsection{Observed quantitative and qualitative indicators}

Table~\ref{tab:validation} summarizes the quantitative indicators recorded across the 22 evaluation profiles. Figure~\ref{fig:eval_chart} displays the performance metrics breakdown.

\begin{table}[H]
\caption{System performance metrics across 22 evaluation test profiles}
\label{tab:validation}
\centering
\begin{tabularx}{\textwidth}{X r}
\toprule
\textbf{Evaluation indicator} & \textbf{Observed value} \\
\midrule
Total benchmark customer profiles & 22 \\
Profiles successfully yielding a complete routine & 21 (95.45\%) \\
Profiles returning zero results (catalogue gap) & 1 (4.55\%) \\
Critical technical or runtime failures & \textbf{0 (0.00\%)} \\
Factual product or price hallucinations & \textbf{0 (0.00\%)} \\
Average response latency (FastAPI deterministic) & 38\,ms \\
Average response latency (Conversational agent with LLM) & 1.84\,s \\
\bottomrule
\end{tabularx}
\end{table}

\placeholderscreenshot{evaluation_benchmark_chart.png}{Evaluation benchmark: routine coherence, budget compliance, and latency metrics across test scenarios}{fig:eval_chart}

\subsection{Observed limitations and catalogue failure modes}

The evaluation successfully isolated true software bugs from catalogue inventory gaps. Table~\ref{tab:gaps} outlines commercial inventory shortages identified through the evaluation suite.

\begin{table}[H]
\caption{Commercial catalogue gaps exposed by scenario evaluation}
\label{tab:gaps}
\centering
\begin{tabularx}{\textwidth}{X r}
\toprule
\textbf{Customer query scenario} & \textbf{Available matching products in catalogue} \\
\midrule
Dandruff accompanied by oily scalp & \textbf{0} \\
Normal skin balanced skincare routine & 2 \\
Damaged and brittle hair repair routine & 2 \\
Dry body skin hydration & 3 \\
\bottomrule
\end{tabularx}
\end{table}

Identifying that the store lacks targeted products for oily dandruff is a valuable commercial insight for the client's procurement team rather than an algorithmic flaw.

\subsection{Positioning against industrial commercial agents}

Table~\ref{tab:industrial_positioning} benchmarks the implemented prototype against leading industrial e-commerce conversational solutions.

\begin{table}[H]
\caption{Positioning of prototype against standard industrial solutions}
\label{tab:industrial_positioning}
\centering
\begin{tabularx}{\textwidth}{L{3.2cm} X X}
\toprule
\textbf{Evaluation dimension} & \textbf{Industrial commercial solutions} & \textbf{Our implemented system} \\
\midrule
Product grounding & Vector search over text chunks; vulnerable to hallucinations & Strict structural tool-calling; mathematically impossible to invent items \\
Attribute dependency & Requires clean, pre-tagged catalogue attributes & Fully automated regex and IDF pipeline extracting metadata from raw prose \\
Operating cost & Monthly SaaS fees + costly per-token inference & Zero marginal cost deterministic engine; low-cost LLM conversational layer \\
Explainability & Opaque black-box neural ranking & Transparent, human-auditable scoring weights and matching justification \\
Routine synthesis & Returns flat, unordered lists of related items & Assembles coherent, multi-step sequential skincare routines \\
\bottomrule
\end{tabularx}
\end{table}

\subsection{Synthesis of evaluation results}

The empirical evaluation proves that the hybrid approach successfully resolves the conflict between generative flexibility and e-commerce reliability. The deterministic scoring core guarantees zero hallucinations, accurate pricing, and rapid response times, while the ReAct agent offers an intuitive, empathetic conversational experience.

\section{Conclusion}

This chapter presented the implementation, containerisation, and systematic evaluation of the recommendation system. By decoupling deterministic scoring from generative dialogue, encapsulating microservices via Docker, and applying rigorous empirical validation across 22 scenarios, the project achieved a robust, explainable, and cost-efficient solution tailored to the client's e-commerce needs.

\cleardoublepage
